\appendix

\section{Additional Experiments}
\label{sec:appendix_experiments}

This appendix reports complementary evidence supporting the main-paper claims and clarifies controls that disentangle interface effects from optimization artifacts.
We report a vocabulary-budget sweep that traces interface trade-offs.

\subsection{Appendix: Vocabulary-Budget Frontier}
\label{sec:appendix_frontier}

Finally, we sweep vocabulary budgets to trace empirical rate--distortion trade-offs and tokenization-time behavior under deployment-aligned execution.
For each vocabulary size, we train each tokenizer on the same training serialization and evaluate (avg./P95) token length, surrogate distortion $\widehat{\Delta}(\tau)$, and tokenization time.
This sweep illustrates how interface constraints and distortion-aware induction shift the frontier, yielding better operating points under tight budgets.

Table~\ref{tab:appendix_frontier} reports the aggregated sweep results.
End-to-end accuracy frontiers require full encoder training at each budget; here we use the sweep to characterize the interface-level operating points (rate, distortion, and runtime).
\begin{table}[h]
\centering
\scriptsize
\setlength{\tabcolsep}{3.5pt}
\begin{tabular*}{\columnwidth}{@{\extracolsep{\fill}}rlccccc@{}}
\toprule
Vocab & Tokenizer & avg\_len & p95\_len & Distortion $\hat{\Delta}$ & tokenize ms/sample \\
\midrule
256 & CIT & 120.6 & 132 & 0.258 & \textbf{0.03} \\
256 & BPE & \textbf{116.4} & \textbf{128} & \textbf{0.187} & 0.05 \\
256 & WordPiece & 124.5 & 136 & 0.202 & 0.05 \\
\midrule
512 & CIT & \textbf{99.6} & \textbf{111} & 0.206 & \textbf{0.03} \\
512 & BPE & 107.6 & 113 & 0.186 & 0.04 \\
512 & WordPiece & 108.1 & 114 & \textbf{0.178} & 0.05 \\
\midrule
1024 & CIT & \textbf{72.2} & \textbf{82} & \textbf{0.144} & \textbf{0.03} \\
1024 & BPE & 106.8 & 112 & 0.203 & 0.04 \\
1024 & WordPiece & 107.0 & 112 & 0.199 & 0.04 \\
\midrule
2048 & CIT & \textbf{60.5} & \textbf{69} & \textbf{0.138} & \textbf{0.03} \\
2048 & BPE & 106.3 & 111 & 0.206 & 0.09 \\
2048 & WordPiece & 106.4 & 112 & 0.186 & 0.04 \\
\midrule
4096 & CIT & \textbf{58.0} & \textbf{60} & \textbf{0.123} & \textbf{0.03} \\
4096 & BPE & 106.1 & 111 & 0.207 & 0.04 \\
4096 & WordPiece & 106.1 & 111 & 0.189 & 0.04 \\
\bottomrule
\end{tabular*}
\caption{Appendix: vocab-budget frontier sweep. Vocabulary sizes are shown in increasing order. All configurations achieve the same accuracy (0.761); thus only efficiency and distortion metrics are reported. Within each vocabulary budget, tokenizers are ordered with CIT first. Best values (lower is better) are shown in bold.}
\label{tab:appendix_frontier}
\end{table}

\paragraph{Note.}
For readability, we present this sweep as a table and omit the corresponding plots.

\paragraph{Reproducibility.}
All main-paper and appendix experiments are implemented as deterministic scripts with fixed serialization and compute-matched training budgets, and the full suite can be run end-to-end from a single driver that also exports publication-ready tables and figures.

\section{Mathematical Proofs}
\label{app:proofs}

This appendix collects deferred proofs that support interface-level claims in the main text.

\subsection{Non-negativity of directed semantic distortion}
\label{app:distortion_nonneg}

\begin{proposition}[Non-negativity of directed semantic distortion]
\label{prop:distortion_nonneg}
For any deterministic, prefix-emitting tokenizer $\tau$ and any prefix length $t$,
\[
\mathcal{R}_\tau(t) \;\ge\; \mathcal{R}^\star(t),
\]
where $\mathcal{R}^\star(t)$ and $\mathcal{R}_\tau(t)$ are defined in
Eq.~\eqref{eq:bayes_risk_raw} and Eq.~\eqref{eq:bayes_risk_tok}, respectively.
Consequently, the directed semantic distortion $\Delta(\tau)$ in Eq.~\eqref{eq:semantic_distortion} is non-negative.
\end{proposition}

\begin{proof}
Fix $t$ and let $Z^{(t)}=\tau(X_{1:t})$.
Any token-observing predictor $q(\,\cdot \mid Z^{(t)})$ induces a raw-prefix predictor
$\tilde{q}(\,\cdot \mid X_{1:t}) \triangleq q(\,\cdot \mid \tau(X_{1:t}))$ by composition with $\tau$.
Since $\mathcal{R}^\star(t)$ is the minimum achievable risk when observing $X_{1:t}$,
\[
\mathcal{R}^\star(t)
\;\le\;
\mathbb{E}\!\left[\ell\!\left(\tilde{q}(\,\cdot \mid X_{1:t}),Y\right)\right]
\;=\;
\mathbb{E}\!\left[\ell\!\left(q(\,\cdot \mid Z^{(t)}),Y\right)\right]
\]
for every such $q$.
Taking the minimum over $q$ yields $\mathcal{R}^\star(t)\le \mathcal{R}_\tau(t)$.
Averaging over $t$ and taking expectations gives $\Delta(\tau)\ge 0$.
\end{proof}

\subsection{Proof of Proposition~\ref{prop:irrecoverability}}
\label{app:proof_irrecoverability}

\begin{proof}[Proof of Proposition~\ref{prop:irrecoverability}]
Let $X_{1:t}$ and $X'_{1:t}$ satisfy $\tau(X_{1:t})=\tau(X'_{1:t})=Z^{(t)}$.
Because $\tau$ is deterministic and prefix-emitting, the tokenized prefix is a deterministic function of the raw prefix.
Therefore, any token-observing predictor $q(\,\cdot \mid Z^{(t)})$ must produce the same conditional predictive distribution for both raw prefixes, since its input is identical.

To show that this can induce strictly positive excess Bayes risk, consider any strictly proper scoring rule $\ell$ (e.g., log loss / cross-entropy), for which the Bayes-optimal prediction is unique and equals the true conditional distribution.
If both prefixes occur with nonzero probability and their Bayes-optimal semantic states differ, i.e.,
$P^\star(\,\cdot \mid X_{1:t}) \neq P^\star(\,\cdot \mid X'_{1:t})$,
then conditioning on $Z^{(t)}$ forces the Bayes-optimal token-observing predictor to predict the mixture
$P^\star(\,\cdot \mid Z^{(t)})$, which is strictly worse (in expected loss) than predicting each conditional distribution separately.
Equivalently, the Bayes risk strictly decreases under the finer observation $X_{1:t}$ compared to the coarser observation $Z^{(t)}$, yielding
$\mathcal{R}_\tau(t) > \mathcal{R}^\star(t)$.
This strict gap is induced by information discarded by the interface and therefore cannot be eliminated by modifying the downstream model without changing $\tau$.
\end{proof}

\subsection{Proof of Proposition~\ref{prop:greedy_descent}}
\label{app:proof_greedy_descent}

\begin{proof}[Proof of Proposition~\ref{prop:greedy_descent}]
By definition of compression gain (Eq.~\eqref{eq:gain_def}),
\[
R(\tau_{\mathcal{V}_{k+1}})
=
R(\tau_{\mathcal{V}_k}) - G(c_k;\mathcal{V}_k).
\]
Under the greedy selection rule (Eq.~\eqref{eq:alg_selection_gain}), $c_k$ maximizes $G(c;\mathcal{V}_k)$ over
contract-feasible candidates.
With the natural convention that the procedure may perform a no-op when all candidates have non-positive gain (equivalently, include a null candidate with gain $0$), we have $G(c_k;\mathcal{V}_k)\ge 0$ for every $k$.
Thus $R(\tau_{\mathcal{V}_{k+1}})\le R(\tau_{\mathcal{V}_k})$ for all $k$.
Since $\widehat{\mathcal{J}}_\lambda$ coincides with $R$ on contract-feasible tokenizers (Eq.~\eqref{eq:relaxed_objective}), the same inequality holds for $\widehat{\mathcal{J}}_\lambda$.
\end{proof}
